% GENERATED FILE. Edit canonical lessons, then run npm run generate:books.
% canonical-source-hash: fnv1a64:feacc937155c695f
% canonical-lessons: SA-C252-valli, SA-C252-kantakah, SA-C252-kalika, SA-C252-mallika, SA-C252-tulasi

\chapter{Vine, Thorn, Bud, Jasmine, Holy basil}
\label{ch:sa-vine-thorn-bud-jasmine-holy-basil}
\hlchaptermodality{252}

\begin{chapteropening}
\textbf{By the end of this chapter:} \emph{I can say a vine, a thorn, a bud, jasmine and holy basil.}

Complete the last lesson of chapter 252: I can say a vine, a thorn, a bud, jasmine and holy basil.
\end{chapteropening}

\section[\texorpdfstring{\sk{वल्ली} (vallī)}{वल्ली (vallī)}]{\sk{वल्ली} (vallī) — a vine}

\label{lesson:SA-C252-valli}



\begin{quote}
Before the new one: say the Sanskrit for nature, then the Sanskrit for mud.
\end{quote}

\subsection*{You'll want to know: \sk{वल्ली}}
\textbf{\sk{वल्ली}} — \emph{vallī} — \textquotedblleft{}a vine\textquotedblright{}.

\begin{grammarlens}[title={What you've built}]
One more word for this chapter.
\end{grammarlens}

\subsection*{Your turn}
\begin{itemize}
\raggedright
  \item \emph{Say it:} \emph{vallī}
  \item \emph{Say it:} \emph{vallī}, once more
  \item \emph{Say it:} say \emph{vallī}
  \item \emph{Recall:} say the Sanskrit for nature, then the Sanskrit for mud, then say \emph{vallī} again
\end{itemize}

\begin{tcolorbox}[breakable,colback=teal!4,colframe=teal!35!black,title={Before you move on}]
What does \sk{वल्ली} mean? (\textquotedblleft{}A vine\textquotedblright{}.) Say it once more.
\end{tcolorbox}
\section[\texorpdfstring{\sk{कण्टकः} (kaṇṭakaḥ)}{कण्टकः (kaṇṭakaḥ)}]{\sk{कण्टकः} (kaṇṭakaḥ) — a thorn}

\label{lesson:SA-C252-kantakah}



\begin{quote}
Before the new one: say the Sanskrit for mud, then the Sanskrit for a vine.
\end{quote}

\subsection*{You'll want to know: \sk{कण्टकः}}
\textbf{\sk{कण्टकः}} — \emph{kaṇṭakaḥ} — \textquotedblleft{}a thorn\textquotedblright{}.

\begin{grammarlens}[title={What you've built}]
One more word for this chapter.
\end{grammarlens}

\subsection*{Your turn}
\begin{itemize}
\raggedright
  \item \emph{Say it:} \emph{kaṇṭakaḥ}
  \item \emph{Say it:} \emph{kaṇṭakaḥ}, once more
  \item \emph{Say it:} say \emph{kaṇṭakaḥ}
  \item \emph{Recall:} say the Sanskrit for mud, then the Sanskrit for a vine, then say \emph{kaṇṭakaḥ} again
\end{itemize}

\begin{tcolorbox}[breakable,colback=teal!4,colframe=teal!35!black,title={Before you move on}]
What does \sk{कण्टकः} mean? (\textquotedblleft{}A thorn\textquotedblright{}.) Say it once more.
\end{tcolorbox}
\section[\texorpdfstring{\sk{कलिका} (kalikā)}{कलिका (kalikā)}]{\sk{कलिका} (kalikā) — a bud}

\label{lesson:SA-C252-kalika}



\begin{quote}
Before the new one: say the Sanskrit for a vine, then the Sanskrit for a thorn.
\end{quote}

\subsection*{You'll want to know: \sk{कलिका}}
\textbf{\sk{कलिका}} — \emph{kalikā} — \textquotedblleft{}a bud\textquotedblright{}.

\begin{grammarlens}[title={What you've built}]
One more word for this chapter.
\end{grammarlens}

\subsection*{Your turn}
\begin{itemize}
\raggedright
  \item \emph{Say it:} \emph{kalikā}
  \item \emph{Say it:} \emph{kalikā}, once more
  \item \emph{Say it:} say \emph{kalikā}
  \item \emph{Recall:} say the Sanskrit for a vine, then the Sanskrit for a thorn, then say \emph{kalikā} again
\end{itemize}

\begin{tcolorbox}[breakable,colback=teal!4,colframe=teal!35!black,title={Before you move on}]
What does \sk{कलिका} mean? (\textquotedblleft{}A bud\textquotedblright{}.) Say it once more.
\end{tcolorbox}
\section[\texorpdfstring{\sk{मल्लिका} (mallikā)}{मल्लिका (mallikā)}]{\sk{मल्लिका} (mallikā) — jasmine}

\label{lesson:SA-C252-mallika}



\begin{quote}
Before the new one: say the Sanskrit for a thorn, then the Sanskrit for a bud.
\end{quote}

\subsection*{You'll want to know: \sk{मल्लिका}}
\textbf{\sk{मल्लिका}} — \emph{mallikā} — \textquotedblleft{}jasmine\textquotedblright{}.

\begin{grammarlens}[title={What you've built}]
One more word for this chapter.
\end{grammarlens}

\subsection*{Your turn}
\begin{itemize}
\raggedright
  \item \emph{Say it:} \emph{mallikā}
  \item \emph{Say it:} \emph{mallikā}, once more
  \item \emph{Say it:} say \emph{mallikā}
  \item \emph{Recall:} say the Sanskrit for a thorn, then the Sanskrit for a bud, then say \emph{mallikā} again
\end{itemize}

\begin{tcolorbox}[breakable,colback=teal!4,colframe=teal!35!black,title={Before you move on}]
What does \sk{मल्लिका} mean? (\textquotedblleft{}Jasmine\textquotedblright{}.) Say it once more.
\end{tcolorbox}
\section[\texorpdfstring{\sk{तुलसी} (tulasī)}{तुलसी (tulasī)}]{\sk{तुलसी} (tulasī) — holy basil}

\label{lesson:SA-C252-tulasi}



\begin{quote}
Before the new one: say the Sanskrit for a bud, then the Sanskrit for jasmine.
\end{quote}

\subsection*{You'll want to know: \sk{तुलसी}}
\textbf{\sk{तुलसी}} — \emph{tulasī} — \textquotedblleft{}holy basil\textquotedblright{}.

\begin{grammarlens}[title={What you've built}]
One more word for this chapter.
\end{grammarlens}

\subsection*{Your turn}
\begin{itemize}
\raggedright
  \item \emph{Say it:} \emph{tulasī}
  \item \emph{Say it:} \emph{tulasī}, once more
  \item \emph{Say it:} say \emph{tulasī}
  \item \emph{Recall:} say the Sanskrit for a bud, then the Sanskrit for jasmine, then say \emph{tulasī} again
\end{itemize}

\begin{tcolorbox}[breakable,colback=teal!4,colframe=teal!35!black,title={Before you move on}]
What does \sk{तुलसी} mean? (\textquotedblleft{}Holy basil\textquotedblright{}.) Say it once more.
\end{tcolorbox}
